% Official Solution (Problem Sheet 13, Exercise 1)
\begin{exercisesolution}[Solution to \cref{exc:explicit_inversion_matrices}]
\label{sol:explicit_inversion_matrices}
\begin{enumerate}[label=\textbf{(\alph*)}]
    \item \textbf{Matrix $A = \begin{pmatrix} 1 & 1 & 1 \\ 1 & 1 & 1 \\ 1 & 1 & 1 \end{pmatrix}$:}
    Applying the row operations $R_2 - R_1 \to R_2$ and $R_3 - R_1 \to R_3$ yields:
    \[
    \begin{pmatrix} 1 & 1 & 1 \\ 0 & 0 & 0 \\ 0 & 0 & 0 \end{pmatrix}.
    \]
    Since $\rank(A) = 1 < 3$, the matrix $A$ is \textbf{not invertible}.

    \item \textbf{Matrix $B = \begin{pmatrix} 1 & 0 & 0 \\ 1 & 1 & 0 \\ 1 & 1 & 1 \end{pmatrix}$:}
    Form the augmented matrix $(B \mid I_3)$ and perform Gaussian elimination:
    \begin{align*}
    \left(\begin{array}{ccc|ccc}
    1 & 0 & 0 & 1 & 0 & 0 \\
    1 & 1 & 0 & 0 & 1 & 0 \\
    1 & 1 & 1 & 0 & 0 & 1
    \end{array}\right)
    &\xrightarrow{\substack{R_2 - R_1 \to R_2 \\ R_3 - R_1 \to R_3}}
    \left(\begin{array}{ccc|ccc}
    1 & 0 & 0 & 1 & 0 & 0 \\
    0 & 1 & 0 & -1 & 1 & 0 \\
    0 & 1 & 1 & -1 & 0 & 1
    \end{array}\right) \\
    &\xrightarrow{R_3 - R_2 \to R_3}
    \left(\begin{array}{ccc|ccc}
    1 & 0 & 0 & 1 & 0 & 0 \\
    0 & 1 & 0 & -1 & 1 & 0 \\
    0 & 0 & 1 & 0 & -1 & 1
    \end{array}\right).
    \end{align*}
    Thus $B$ is invertible with inverse:
    \[
    B^{-1} = \begin{pmatrix} 1 & 0 & 0 \\ -1 & 1 & 0 \\ 0 & -1 & 1 \end{pmatrix}.
    \]

    \item \textbf{Matrix $C = \begin{pmatrix} 2 & 1 & 1 & 1 \\ 1 & 2 & 1 & 1 \\ 1 & 1 & 2 & 1 \\ 1 & 1 & 1 & 2 \end{pmatrix}$:}
    Notice $C = I_4 + J_4$ where $J_4$ is the $4 \times 4$ all-ones matrix.
    Sum the other three rows into the first, one operation at a time, by
    $R_1 + R_2 \to R_1$, then $R_1 + R_3 \to R_1$, then $R_1 + R_4 \to R_1$. Since every
    column of $C$ sums to $5$, the first row becomes
    \[
    (5, 5, 5, 5 \mid 1, 1, 1, 1).
    \]
    Now $\frac{1}{5} \cdot R_1 \to R_1$ gives $(1, 1, 1, 1 \mid \frac{1}{5}, \frac{1}{5}, \frac{1}{5}, \frac{1}{5})$.
    Subtracting this new first row from each of the others, by $R_i - R_1 \to R_i$ for $i = 2, 3, 4$, gives $e_i\transp \mid -\frac{1}{5}, \dots, \frac{4}{5}, \dots$.
    Finally $R_1 - R_2 \to R_1$, then $R_1 - R_3 \to R_1$, then $R_1 - R_4 \to R_1$ leaves $I_4$ on the left.
    Thus $C$ is invertible with inverse:
    \[
    C^{-1} = \frac{1}{5} \begin{pmatrix}
    4 & -1 & -1 & -1 \\
    -1 & 4 & -1 & -1 \\
    -1 & -1 & 4 & -1 \\
    -1 & -1 & -1 & 4
    \end{pmatrix}.
    \]

    \item \textbf{Matrix $D = \begin{pmatrix} 1 & 2 & -3 & 1 \\ -1 & 3 & -3 & -2 \\ 2 & 0 & 1 & 5 \\ 3 & 1 & -2 & 6 \end{pmatrix}$:}
    Applying row elimination to $(D \mid I_4)$ with compact notation:
    \begin{align*}
    (D \mid I_4) &=
    \left(\begin{array}{cccc|cccc}
    1 & 2 & -3 & 1 & 1 & 0 & 0 & 0 \\
    -1 & 3 & -3 & -2 & 0 & 1 & 0 & 0 \\
    2 & 0 & 1 & 5 & 0 & 0 & 1 & 0 \\
    3 & 1 & -2 & 6 & 0 & 0 & 0 & 1
    \end{array}\right) \\
    &\xrightarrow{\substack{R_2 + R_1 \to R_2 \\ R_3 - 2R_1 \to R_3 \\ R_4 - 3R_1 \to R_4}}
    \left(\begin{array}{cccc|cccc}
    1 & 2 & -3 & 1 & 1 & 0 & 0 & 0 \\
    0 & 5 & -6 & -1 & 1 & 1 & 0 & 0 \\
    0 & -4 & 7 & 3 & -2 & 0 & 1 & 0 \\
    0 & -5 & 7 & 3 & -3 & 0 & 0 & 1
    \end{array}\right) \\
    &\xrightarrow{\substack{R_4 + R_2 \to R_4 \\ R_3 + \frac{4}{5}R_2 \to R_3}}
    \left(\begin{array}{cccc|cccc}
    1 & 2 & -3 & 1 & 1 & 0 & 0 & 0 \\
    0 & 5 & -6 & -1 & 1 & 1 & 0 & 0 \\
    0 & 0 & \frac{11}{5} & \frac{11}{5} & -\frac{6}{5} & \frac{4}{5} & 1 & 0 \\
    0 & 0 & 1 & 2 & -2 & 1 & 0 & 1
    \end{array}\right) \\
    &\xrightarrow{\frac{5}{11}R_3 \to R_3}
    \left(\begin{array}{cccc|cccc}
    1 & 2 & -3 & 1 & 1 & 0 & 0 & 0 \\
    0 & 5 & -6 & -1 & 1 & 1 & 0 & 0 \\
    0 & 0 & 1 & 1 & -\frac{6}{11} & \frac{4}{11} & \frac{5}{11} & 0 \\
    0 & 0 & 1 & 2 & -2 & 1 & 0 & 1
    \end{array}\right) \\
    &\xrightarrow{R_4 - R_3 \to R_4}
    \left(\begin{array}{cccc|cccc}
    1 & 2 & -3 & 1 & 1 & 0 & 0 & 0 \\
    0 & 5 & -6 & -1 & 1 & 1 & 0 & 0 \\
    0 & 0 & 1 & 1 & -\frac{6}{11} & \frac{4}{11} & \frac{5}{11} & 0 \\
    0 & 0 & 0 & 1 & -\frac{16}{11} & \frac{7}{11} & -\frac{5}{11} & 1
    \end{array}\right).
    \end{align*}
    Back-substituting upwards yields the reduced identity matrix $(I_4 \mid D^{-1})$:
    \[
    D^{-1} = \frac{1}{11} \begin{pmatrix}
    35 & -16 & 13 & -22 \\
    11 & 0 & 11 & -11 \\
    10 & -3 & 10 & -11 \\
    -16 & 7 & -5 & 11
    \end{pmatrix}. \qedhere
    \]
\end{enumerate}
\exinfo{Official Solution (Problem Sheet 13, Exercise 1). Explicitly inverts 4 matrices via Gauss-Jordan elimination on $(M \mid I)$.}
\end{exercisesolution}

% Official Solution (Problem Sheet 13, Exercise 5)
\begin{exercisesolution}[Solution to \cref{exc:parametric_gaussian_elimination_invertibility}]
\label{sol:parametric_gaussian_elimination_invertibility}
We apply Gaussian elimination to the parameter-dependent matrix $M_\alpha$:
\[
M_\alpha = \begin{pmatrix}
1 & 3 & -4 & 2 \\
2 & 1 & -2 & 1 \\
3 & -1 & -2 & -2\alpha \\
-6 & 2 & 1 & \alpha^2
\end{pmatrix}
\xrightarrow{\substack{R_2 - 2R_1 \to R_2 \\ R_3 - 3R_1 \to R_3 \\ R_4 + 6R_1 \to R_4}}
\begin{pmatrix}
1 & 3 & -4 & 2 \\
0 & -5 & 6 & -3 \\
0 & -10 & 10 & -2\alpha - 6 \\
0 & 20 & -23 & \alpha^2 + 12
\end{pmatrix}.
\]
Next, eliminate the second column below the second pivot:
\[
\xrightarrow{\substack{R_3 - 2R_2 \to R_3 \\ R_4 + 4R_2 \to R_4}}
\begin{pmatrix}
1 & 3 & -4 & 2 \\
0 & -5 & 6 & -3 \\
0 & 0 & -2 & -2\alpha \\
0 & 0 & 1 & \alpha^2
\end{pmatrix}
\xrightarrow{-\frac{1}{2}R_3 \to R_3}
\begin{pmatrix}
1 & 3 & -4 & 2 \\
0 & -5 & 6 & -3 \\
0 & 0 & 1 & \alpha \\
0 & 0 & 1 & \alpha^2
\end{pmatrix}.
\]
Finally, eliminate the third column in the fourth row:
\[
\xrightarrow{R_4 - R_3 \to R_4}
\begin{pmatrix}
1 & 3 & -4 & 2 \\
0 & -5 & 6 & -3 \\
0 & 0 & 1 & \alpha \\
0 & 0 & 0 & \alpha^2 - \alpha
\end{pmatrix}.
\]
The row echelon form has $4$ non-zero pivots if and only if the bottom right entry is non-zero:
\[
\alpha^2 - \alpha \ne 0 \iff \alpha(\alpha - 1) \ne 0 \iff \alpha \notin \{0, 1\}.
\]
Therefore, $M_\alpha$ is invertible if and only if $\alpha \in \mathbb{Q} \setminus \{0, 1\}$.
\exinfo{Official Solution (Problem Sheet 13, Exercise 5). Performs parametric row reduction to determine values of $\alpha$ for matrix invertibility.}
\end{exercisesolution}

% Solution: Gemini / Official Hybrid (Problem Sheet 13, Exercise 6)
\begin{exercisesolution}[Solution to \cref{exc:bruhat_decomposition}]
\label{sol:bruhat_decomposition}
Let $A \in \GL_n(K)$. Let $B \subseteq \GL_n(K)$ denote the subgroup of invertible upper triangular matrices.

\textbf{Proof via Gaussian Elimination:}
For any non-zero vector $v \in K^n$, let $\operatorname{piv}(v) \in \{1, \dots, n\}$ denote the index of the \emph{last} (bottom-most) non-zero entry of $v$.
We perform row reduction using only upper triangular elementary row operations (Type I operations $R_i + \alpha R_j \to R_i$ with $i < j$ and Type III scaling operations $\alpha \cdot R_i \to R_i$).
Left-multiplication by an upper triangular matrix $U \in B$ allows us to clear all entries in row $i$ sitting above any established pivot position $j$ (where $i < j$) without affecting lower rows.

Because $A$ is invertible, no row ever becomes zero. By iteratively applying these upper triangular operations from bottom to top, we can ensure that each row $i$ has its pivot in a unique column $\sigma(i)$, with no other non-zero entries in column $\sigma(i)$ above row $i$.
The permutation $\sigma \in S_n$ defines a permutation matrix $P \in \GL_n(K)$ with $P_{i, \sigma(i)} = 1$.
The reduced matrix $U A$ then has the property that right-multiplication by an upper triangular matrix $U' \in B$ clears all entries to the right of the pivots, reducing $U A U' = P$.
Multiplying by the inverses $B := U^{-1} \in B$ and $B' := (U')^{-1} \in B$ yields the Bruhat decomposition:
\[
A = B P B'. \qedhere
\]
\exinfo{Hybrid Solution (Gemini 3.7 Flash / Official Problem Sheet 13, Exercise 6). Establishes the Bruhat decomposition $A = B P B'$ for $\GL_n(K)$ using triangular row operations and permutation matrices.}
\end{exercisesolution}

% Official Solution (Problem Sheet 14, Multiple Choice Question 1)
\begin{exercisesolution}[Solution to \cref{exc:parametric_3x3_invertibility_mc}]
\label{sol:parametric_3x3_invertibility_mc}
Let $A = \begin{pmatrix} 1 & x & 1 \\ 3 & 3 & x \\ 0 & 3 & 1 \end{pmatrix} \in \M_{3 \times 3}(\mathbb{R})$.
We compute the determinant of $A$ by expanding along the first column:
\begin{align*}
\det(A) &= 1 \cdot \det\begin{pmatrix} 3 & x \\ 3 & 1 \end{pmatrix} - 3 \cdot \det\begin{pmatrix} x & 1 \\ 3 & 1 \end{pmatrix} + 0 \\
&= (3 - 3x) - 3(x - 3) = 3 - 3x - 3x + 9 = 12 - 6x.
\end{align*}
A square matrix is not invertible if and only if its determinant is zero:
\[
\det(A) = 0 \iff 12 - 6x = 0 \iff 6x = 12 \iff x = 2.
\]
Alternatively, applying the row operation $R_2 - 3R_1 \to R_2$ gives $\begin{pmatrix} 1 & x & 1 \\ 0 & 3-3x & x-3 \\ 0 & 3 & 1 \end{pmatrix}$.
When $x = 2$, row 2 is $(0, -3, -1) = -R_3$, producing a zero row upon $R_3 + R_2 \to R_3$.
Thus, $A$ is not invertible if and only if $x = 2$.
\exinfo{Official Solution (Problem Sheet 14, Multiple Choice Question 1). Evaluates the non-invertibility condition $\det(A) = 0$ for a $3 \times 3$ parametric matrix.}
\end{exercisesolution}
